\documentclass[12pt]{article}
\usepackage{amsmath, amssymb, geometry, amsthm, mdframed, xcolor, enumitem}
\geometry{letterpaper, margin=1in}

\title{\textbf{Homework 1: Sets, Partitions, and Cartesian Products}}
\author{Tashfeen Omran}
\date{September 2026}

\begin{document}

\maketitle
\tableofcontents
\newpage

\section{Part 1: Comprehensive Introduction, Context, and Prerequisites}

\subsection{Core Concepts}
Set theory is the fundamental language of modern mathematics. While basic set operations (union, intersection) are intuitive, advancing to higher mathematics requires extending these concepts to collections of sets. 

\begin{itemize}
    \item \textbf{Indexed Collections of Sets:} Instead of looking at just two sets $A$ and $B$, we often deal with an infinite family of sets. We use an \textit{index set} $I$. For every $\alpha \in I$, there is a corresponding set $A_\alpha$. The collection is denoted $\{A_\alpha\}_{\alpha \in I}$. 
    \item \textbf{Partitions:} A partition of a set $A$ is a way of breaking $A$ into mutually exclusive, exhaustive pieces. Formally, a partition is a set of subsets of $A$ that are non-empty, pairwise disjoint (no two share an element), and whose union equals $A$.
    \item \textbf{Cartesian Products:} For two sets $A$ and $B$, the Cartesian product $A \times B$ is the set of all ordered pairs $(a, b)$ where $a \in A$ and $b \in B$. This concept is the bedrock for defining relations, functions, and multidimensional coordinate spaces.
\end{itemize}

\subsection{Intuition and Motivation}
Historically, mathematics was somewhat informal until the late 19th century when Georg Cantor formalized Set Theory. Mathematicians needed a rigorous way to handle "infinity" and collections of objects. For example, understanding continuous functions or the real number line requires talking about the intersection of an infinite number of infinitesimally small intervals. 

Partitions provide the mathematical framework for "classification." If you want to classify integers as even or odd, you are mathematically partitioning the set of integers into two subsets. Cartesian products give us geometry: the $xy$-plane is simply the Cartesian product $\mathbb{R} \times \mathbb{R}$.

\subsection{Key Formulas and Notation}
\begin{itemize}
    \item \textbf{Union of a Collection:} $\bigcup_{\alpha \in I} A_\alpha = \{x \mid x \in A_\alpha \text{ for at least one } \alpha \in I\}$.
    \item \textbf{Intersection of a Collection:} $\bigcap_{\alpha \in I} A_\alpha = \{x \mid x \in A_\alpha \text{ for every } \alpha \in I\}$.
    \item \textbf{Cartesian Product:} $A \times B = \{(a, b) \mid a \in A \text{ and } b \in B\}$.
    \item \textbf{Cardinality of Cartesian Product:} For finite sets, $|A \times B| = |A| \cdot |B|$.
    \item \textbf{Power Set:} $\mathcal{P}(A)$ is the set of all subsets of $A$. $|\mathcal{P}(A)| = 2^{|A|}$.
\end{itemize}

\subsection{Prerequisites}
To master this section, you must be comfortable with:
\begin{enumerate}
    \item \textbf{Interval Notation:} Understanding open $(a, b)$, closed $[a, b]$, and half-open $[a, b)$ intervals of real numbers.
    \item \textbf{Basic Set Operations:} Subsets ($\subseteq$), unions ($\cup$), intersections ($\cap$), and set difference ($\setminus$ or $-$).
    \item \textbf{Limit Intuition:} Basic understanding of what happens to expressions like $1/n$ as $n \to \infty$.
\end{enumerate}

\newpage
\section{Part 2: Detailed Homework Solutions}

\subsection{Section 1.4: Indexed Collections of Sets}

\textbf{1.36. For a real number $r$, define $S_r$ to be the interval $[r-1, r+2]$. Let $A = \{1, 3, 4\}$. Determine $\bigcup_{\alpha \in A} S_\alpha$ and $\bigcap_{\alpha \in A} S_\alpha$.}

\textbf{Solution:}
First, substitute each element of the index set $A$ to find the individual sets:
\begin{align*}
    S_1 &= [1-1, 1+2] = [0, 3] \\
    S_3 &= [3-1, 3+2] = [2, 5] \\
    S_4 &= [4-1, 4+2] = [3, 6]
\end{align*}
The union is the combination of all these intervals:
\[ \bigcup_{\alpha \in A} S_\alpha = [0, 3] \cup [2, 5] \cup [3, 6] = [0, 6] \]
The intersection contains only the elements present in \textit{all three} intervals:
\[ \bigcap_{\alpha \in A} S_\alpha = [0, 3] \cap [2, 5] \cap [3, 6] \]
Notice that the only number present in $[0, 3]$ and $[3, 6]$ is $3$. Since $3$ is also in $[2, 5]$, the intersection is just the single element:
\[ \bigcap_{\alpha \in A} S_\alpha = \{3\} \]

\hrulefill

\textbf{1.37. Let $A = \{1, 2, 5\}$, $B = \{0, 2, 4\}$, $C = \{2, 3, 4\}$ and $S = \{A, B, C\}$. Determine $\bigcup_{X \in S} X$ and $\bigcap_{X \in S} X$.}

\textbf{Solution:}
The collection $S$ contains three sets. 
The union is all distinct elements present in any of the sets:
\[ \bigcup_{X \in S} X = A \cup B \cup C = \{0, 1, 2, 3, 4, 5\} \]
The intersection is the elements present in every set. The only element common to $A$, $B$, and $C$ is $2$:
\[ \bigcap_{X \in S} X = A \cap B \cap C = \{2\} \]

\hrulefill

\textbf{1.38. For a real number $r$, define $A_r = \{r^2\}$, $B_r$ as the closed interval $[r-1, r+1]$, and $C_r$ as the interval $(r, \infty)$. For $S = \{1, 2, 4\}$, determine:}
(a) $\bigcup_{\alpha \in S} A_\alpha$ and $\bigcap_{\alpha \in S} A_\alpha$
(b) $\bigcup_{\alpha \in S} B_\alpha$ and $\bigcap_{\alpha \in S} B_\alpha$
(c) $\bigcup_{\alpha \in S} C_\alpha$ and $\bigcap_{\alpha \in S} C_\alpha$

\textbf{Solution:}
\textbf{Part (a):} $A_1 = \{1\}$, $A_2 = \{4\}$, $A_4 = \{16\}$.
\[ \bigcup_{\alpha \in S} A_\alpha = \{1, 4, 16\}, \quad \bigcap_{\alpha \in S} A_\alpha = \emptyset \]
\textbf{Part (b):} $B_1 = [0, 2]$, $B_2 = [1, 3]$, $B_4 = [3, 5]$.
\[ \bigcup_{\alpha \in S} B_\alpha = [0, 2] \cup [1, 3] \cup [3, 5] = [0, 5] \]
\[ \bigcap_{\alpha \in S} B_\alpha = [0, 2] \cap [1, 3] \cap [3, 5] = \emptyset \]
\textbf{Part (c):} $C_1 = (1, \infty)$, $C_2 = (2, \infty)$, $C_4 = (4, \infty)$.
Because these intervals are strictly nested ($C_4 \subset C_2 \subset C_1$), their union is the largest set and the intersection is the smallest set:
\[ \bigcup_{\alpha \in S} C_\alpha = (1, \infty), \quad \bigcap_{\alpha \in S} C_\alpha = (4, \infty) \]

\hrulefill

\textbf{1.39. Let $A = \{a, b, \ldots, z\}$ be the set of letters of the alphabet. For $\alpha \in A$, let $A_\alpha$ consist of $\alpha$ and the two letters that follow it, where $A_y = \{y, z, a\}$ and $A_z = \{z, a, b\}$. Find a set $S \subseteq A$ of smallest cardinality such that $\bigcup_{\alpha \in S} A_\alpha = A$. Explain.}

\textbf{Solution:}
Each set $A_\alpha$ covers exactly 3 letters. The alphabet contains 26 letters. To cover all 26 letters using sets of size 3, we must select at least $\lceil 26 / 3 \rceil = 9$ sets.
To minimize overlap, we pick letters spaced exactly 3 positions apart. Let's start at $a$:
\begin{align*}
    A_a &= \{a, b, c\} \\
    A_d &= \{d, e, f\} \\
    A_g &= \{g, h, i\} \\
    A_j &= \{j, k, l\} \\
    A_m &= \{m, n, o\} \\
    A_p &= \{p, q, r\} \\
    A_s &= \{s, t, u\} \\
    A_v &= \{v, w, x\} \\
    A_y &= \{y, z, a\}
\end{align*}
The union of these 9 sets covers all 26 letters. Thus, the index set is:
\[ S = \{a, d, g, j, m, p, s, v, y\} \]
The smallest cardinality is $9$.

\hrulefill

\textbf{1.40. For $i \in \mathbb{Z}$, let $A_i = \{i-1, i+1\}$. Determine the following:}
(a) $\bigcup_{i=1}^5 A_{2i}$ \quad (b) $\bigcup_{i=1}^5 (A_i \cap A_{i+1})$ \quad (c) $\bigcup_{i=1}^5 (A_{2i-1} \cap A_{2i+1})$

\textbf{Solution:}
\textbf{Part (a):} We evaluate $A_{2i}$ for $i=1$ to $5$:
$A_2 = \{1, 3\}$, $A_4 = \{3, 5\}$, $A_6 = \{5, 7\}$, $A_8 = \{7, 9\}$, $A_{10} = \{9, 11\}$.
\[ \bigcup_{i=1}^5 A_{2i} = \{1, 3, 5, 7, 9, 11\} \]
\textbf{Part (b):} Evaluate $A_i \cap A_{i+1}$.
$A_i = \{i-1, i+1\}$ and $A_{i+1} = \{i, i+2\}$. Since elements differ by parities (one set has even and the other has odd relative to $i$), their intersection is always empty: $A_i \cap A_{i+1} = \emptyset$.
\[ \bigcup_{i=1}^5 \emptyset = \emptyset \]
\textbf{Part (c):} Evaluate $A_{2i-1} \cap A_{2i+1}$.
$A_{2i-1} = \{2i-2, 2i\}$ and $A_{2i+1} = \{2i, 2i+2\}$. The shared element is exactly $2i$.
So, $\bigcup_{i=1}^5 \{2i\} = \{2, 4, 6, 8, 10\}$.

\hrulefill

\textbf{1.41. Find an indexed collection $\{A_n\}_{n \in \mathbb{N}}$ of distinct sets satisfying the given conditions.}
(a) $\bigcap_{n=1}^\infty A_n = \{0\}$ and $\bigcup_{n=1}^\infty A_n = [0, 1]$
(b) $\bigcap_{n=1}^\infty A_n = \{-1, 0, 1\}$ and $\bigcup_{n=1}^\infty A_n = \mathbb{Z}$

\textbf{Solution:}
\textbf{Part (a):} We need intervals that shrink down to $0$ but whose largest extent is $[0, 1]$.
Let $A_n = [0, \frac{1}{n}]$. 
Union: The largest interval is at $n=1$, which is $[0, 1]$.
Intersection: As $n \to \infty$, $\frac{1}{n} \to 0$, leaving only $\{0\}$. All sets are distinct.
\textbf{Part (b):} We need sets whose core is $\{-1, 0, 1\}$ and that expand to cover all integers.
Let $A_n = \{-n, \dots, n\}$. However, $A_1 = \{-1, 0, 1\}$. To ensure the intersection is EXACTLY this, we anchor it:
Let $A_n = \{k \in \mathbb{Z} \mid -n \le k \le n\}$.
Intersection: $\bigcap_{n=1}^\infty A_n = A_1 = \{-1, 0, 1\}$.
Union: As $n \to \infty$, it spans all integers, so $\bigcup A_n = \mathbb{Z}$.

\hrulefill

\textbf{1.42. For each of the following collections of sets, define a set $A_n$ for each $n \in \mathbb{N}$ such that the indexed collection $\{A_n\}_{n \in \mathbb{N}}$ is precisely the given collection. Then find both the union and intersection.}
(a) $\{[1, 2+1), [1, 2+\frac{1}{2}), [1, 2+\frac{1}{3}), \dots\}$
(b) $\{(-1, 2), (-\frac{3}{2}, 4), (-\frac{5}{3}, 6), (-\frac{7}{4}, 8), \dots\}$

\textbf{Solution:}
\textbf{Part (a):} By matching terms to $n$, we define:
\[ A_n = \left[1, 2 + \frac{1}{n}\right) \]
Since $2 + \frac{1}{n}$ decreases as $n$ increases, the intervals are strictly shrinking.
Union is the largest set ($n=1$): $\bigcup_{n=1}^\infty A_n = [1, 3)$.
Intersection is the limit as $n \to \infty$. $2 + \frac{1}{n}$ approaches $2$ from above. Since $x < 2+\frac{1}{n}$ for all $n$, $x$ can equal $2$. Thus: $\bigcap_{n=1}^\infty A_n = [1, 2]$.

\textbf{Part (b):} Pattern matching for upper bounds: $2, 4, 6, 8 \implies 2n$.
Pattern for lower bounds: $-1, -1.5, -1.666, -1.75 \implies -2 + \frac{1}{n}$.
\[ A_n = \left(-2 + \frac{1}{n}, 2n\right) \]
These intervals are strictly expanding. 
Intersection is the smallest set ($n=1$): $\bigcap_{n=1}^\infty A_n = (-1, 2)$.
Union is the limit as $n \to \infty$: $\bigcup_{n=1}^\infty A_n = (-2, \infty)$.

\hrulefill

\textbf{1.43. For $r \in \mathbb{R}^+$, let $A_r = \{x \in \mathbb{R} : |x| < r\}$. Determine $\bigcup_{r \in \mathbb{R}^+} A_r$ and $\bigcap_{r \in \mathbb{R}^+} A_r$.}

\textbf{Solution:}
The condition $|x| < r$ defines the open interval $A_r = (-r, r)$.
Since $r$ can be arbitrarily large, the union covers all real numbers:
\[ \bigcup_{r \in \mathbb{R}^+} A_r = \mathbb{R} \]
As $r$ approaches $0$ from the right, the interval $(-r, r)$ shrinks. The only number whose absolute value is strictly less than every positive real number $r$ is $0$.
\[ \bigcap_{r \in \mathbb{R}^+} A_r = \{0\} \]

\hrulefill

\textbf{1.44. Each of the following sets is a subset of $A = \{1, 2, \dots, 10\}$: (List of $A_1$ through $A_{13}$ omitted for brevity). Find a set $I \subseteq \{1, 2, \dots, 13\}$ such that for every two distinct elements $j, k \in I$, $A_j \cap A_k = \emptyset$ and $|\bigcup_{i \in I} A_i|$ is maximum.}

\textbf{Solution:}
We are tasked with finding a collection of mutually disjoint sets from the list whose union contains the maximum possible number of elements. 
After exhaustive pairwise testing of the 13 sets, it becomes clear that no three sets are mutually disjoint. Thus, the index set $I$ can contain at most 2 elements.
We check all pairs to maximize the sum of their cardinalities without overlap.
Consider $A_1 = \{1, 5, 7, 9, 10\}$ (size 5) and $A_4 = \{2, 4, 8\}$ (size 3).
Their intersection is empty: $A_1 \cap A_4 = \emptyset$.
Their union is $\{1, 2, 4, 5, 7, 8, 9, 10\}$, which has $8$ elements. 
No other disjoint pair yields a sum of 9 or 10. Thus, a valid index set is:
\[ I = \{1, 4\} \]

\hrulefill

\textbf{1.45. For $n \in \mathbb{N}$, let $A_n = (-\frac{1}{n}, 2 - \frac{1}{n})$. Determine $\bigcup_{n \in \mathbb{N}} A_n$ and $\bigcap_{n \in \mathbb{N}} A_n$.}

\textbf{Solution:}
Let's analyze the bounds as $n$ grows:
Lower bound: $-\frac{1}{n}$. As $n \to \infty$, $-\frac{1}{n}$ increases towards $0$.
Upper bound: $2 - \frac{1}{n}$. As $n \to \infty$, $2 - \frac{1}{n}$ increases towards $2$.
\textbf{Union:} Takes the smallest lower bound ($n=1 \implies -1$) and the largest upper bound (limit $\implies 2$).
\[ \bigcup_{n=1}^\infty A_n = (-1, 2) \]
\textbf{Intersection:} Needs elements that are in \textit{every} set. 
For the lower bound, $x > -\frac{1}{n}$ for all $n$. The tightest restriction is $n \to \infty$, implying $x \ge 0$.
For the upper bound, $x < 2 - \frac{1}{n}$ for all $n$. The tightest restriction is $n=1$, implying $x < 1$.
\[ \bigcap_{n=1}^\infty A_n = [0, 1) \]

\hrulefill

\textbf{1.46. Determine the following:}
(a) $\bigcup_{n=1}^\infty \left(-\frac{1}{n}, \frac{1}{n}\right)$ and $\bigcap_{n=1}^\infty \left(-\frac{1}{n}, \frac{1}{n}\right)$
(b) $\bigcup_{n=1}^\infty \left[\frac{n-1}{n}, \frac{n+1}{n}\right]$ and $\bigcap_{n=1}^\infty \left[\frac{n-1}{n}, \frac{n+1}{n}\right]$

\textbf{Solution:}
\textbf{Part (a):} These intervals are strictly shrinking as $n$ increases.
Union is the $n=1$ set: $(-1, 1)$.
Intersection is the limit as $n \to \infty$: $\{0\}$.
\textbf{Part (b):} Rewrite bounds as $\left[1 - \frac{1}{n}, 1 + \frac{1}{n}\right]$.
These intervals are also strictly shrinking toward $1$.
Union is the $n=1$ set: $[0, 2]$.
Intersection is the limit as $n \to \infty$: $\{1\}$.

\hrulefill

\textbf{1.47. Determine the following:}
(a) $\bigcup_{n=1}^\infty \left\{ \sin^2 \frac{n\pi}{2} + \cos^2 \frac{n\pi}{2} \right\}$
(b) $\bigcup_{n=1}^\infty \left\{ \sin \frac{n\pi}{2} + \cos \frac{n\pi}{2} \right\}$

\textbf{Solution:}
\textbf{Part (a):} By the Pythagorean identity, $\sin^2 x + \cos^2 x = 1$ for all real $x$. Thus, for every $n$, the set is simply $\{1\}$.
The infinite union of $\{1\}$ with itself is: $\{1\}$.
\textbf{Part (b):} We evaluate the expression for the first few values of $n$:
$n=1: \sin(\frac{\pi}{2}) + \cos(\frac{\pi}{2}) = 1 + 0 = 1$
$n=2: \sin(\pi) + \cos(\pi) = 0 - 1 = -1$
$n=3: \sin(\frac{3\pi}{2}) + \cos(\frac{3\pi}{2}) = -1 + 0 = -1$
$n=4: \sin(2\pi) + \cos(2\pi) = 0 + 1 = 1$
The values oscillate between $1$ and $-1$. Thus, the sets are either $\{1\}$ or $\{-1\}$.
The union is $\{-1, 1\}$.

\subsection{Section 1.5: Partitions of Sets}

\textbf{1.48. Which of the following are partitions of $A = \{a, b, c, d, e, f, g\}$?}
(a) $S_1 = \{\{a, c, e, g\}, \{b, f\}, \{d\}\}$ \quad (b) $S_2 = \{\{a, b, c, d\}, \{e, f\}\}$
(c) $S_3 = \{A\}$ \quad (d) $S_4 = \{\{a\}, \emptyset, \{b, c, d\}, \{e, f, g\}\}$ \quad (e) $S_5 = \{\{a, c, d\}, \{b, g\}, \{e\}, \{b, f\}\}$

\textbf{Solution:}
A partition must be disjoint, non-empty subsets whose union is $A$.
(a) \textbf{Yes.} Disjoint, union is $A$.
(b) \textbf{No.} Missing element $g$.
(c) \textbf{Yes.} Trivial partition.
(d) \textbf{No.} Contains the empty set $\emptyset$.
(e) \textbf{No.} Not pairwise disjoint (element $b$ is in two subsets).

\hrulefill

\textbf{1.49. Which of the following are partitions of $A = \{1, 2, 3, 4, 5\}$?}
(a) $S_1 = \{\{1, 3\}, \{2, 5\}\}$ \quad (b) $S_2 = \{\{1, 2\}, \{3, 4, 5\}\}$
(c) $S_3 = \{\{1, 2\}, \{2, 3\}, \{3, 4\}, \{4, 5\}\}$ \quad (d) $S_4 = A$

\textbf{Solution:}
(a) \textbf{No.} Missing element $4$.
(b) \textbf{Yes.} Disjoint, union is $A$.
(c) \textbf{No.} Not disjoint.
(d) \textbf{No.} $A$ is a set of numbers, but a partition must be a set of \textit{sets}.

\hrulefill

\textbf{1.50. Let $A = \{1, 2, 3, 4, 5, 6\}$. Give an example of a partition $S$ of $A$ such that $|S| = 3$.}

\textbf{Solution:}
We need to divide $A$ into exactly 3 subsets. 
\[ S = \{\{1, 2\}, \{3, 4\}, \{5, 6\}\} \]

\hrulefill

\textbf{1.51. Give an example of a set $A$ with $|A| = 4$ and two disjoint partitions $S_1$ and $S_2$ of $A$ with $|S_1| = |S_2| = 3$.}

\textbf{Solution:}
Let $A = \{1, 2, 3, 4\}$. We need two partitions of size 3 that share no subsets.
$S_1 = \{\{1, 2\}, \{3\}, \{4\}\}$
$S_2 = \{\{3, 4\}, \{1\}, \{2\}\}$
$S_1 \cap S_2 = \emptyset$, so they are disjoint partitions.

\hrulefill

\textbf{1.52. Give an example of a partition of $\mathbb{N}$ into three subsets.}

\textbf{Solution:}
We can partition natural numbers by their remainder when divided by 3:
\[ S = \{\{3k-2 \mid k \in \mathbb{N}\}, \{3k-1 \mid k \in \mathbb{N}\}, \{3k \mid k \in \mathbb{N}\}\} \]

\hrulefill

\textbf{1.53. Give an example of a partition of $\mathbb{Q}$ into three subsets.}

\textbf{Solution:}
Partition rationals into positive, negative, and zero.
\[ S = \{\mathbb{Q}^+, \mathbb{Q}^-, \{0\}\} \]

\hrulefill

\textbf{1.54. Give an example of three sets $A, S_1, S_2$ such that $S_1$ is a partition of $A$, $S_2$ is a partition of $S_1$ and $|S_2| < |S_1| < |A|$.}

\textbf{Solution:}
Let $A = \{1, 2, 3, 4\}$. ($|A| = 4$)
Let $S_1 = \{\{1, 2\}, \{3\}, \{4\}\}$. ($|S_1| = 3$, partitions $A$)
Let $S_2 = \{\{\{1, 2\}, \{3\}\}, \{\{4\}\}\}$. ($|S_2| = 2$, partitions $S_1$)

\hrulefill

\textbf{1.55. Give an example of a partition of $\mathbb{Z}$ into four subsets.}

\textbf{Solution:}
Partition by congruence modulo 4.
$S = \{A_0, A_1, A_2, A_3\}$, where $A_i = \{4k + i \mid k \in \mathbb{Z}\}$.

\hrulefill

\textbf{1.56. Let $A = \{1, 2, \ldots, 12\}$. Give an example of a partition $S$ of $A$ satisfying: (i) $|S|=5$, (ii) there is a subset $T$ of $S$ such that $|T|=4$ and $|\cup_{X \in T} X| = 10$, and (iii) no element $B \in S$ has $|B|=3$.}

\textbf{Solution:}
We need to partition 12 elements into 5 sets (sizes sum to 12).
Condition (ii) says 4 of these sets contain 10 elements. This implies the 5th set has $12 - 10 = 2$ elements.
So we need four numbers that sum to 10. By condition (iii), no size can be 3.
A valid combination of sizes for the 4 sets is $4, 2, 2, 2$ (since $4+2+2+2 = 10$).
So the five sets have sizes $4, 2, 2, 2, 2$.
\[ S = \{\{1, 2, 3, 4\}, \{5, 6\}, \{7, 8\}, \{9, 10\}, \{11, 12\}\} \]

\hrulefill

\textbf{1.57. A set $S$ is partitioned into two subsets. This produces a partition $P_1$ of size 2. Each is then partitioned into two subsets producing $P_2$ of size 4. This continues until $P_6$. What is $|P_6|$?}

\textbf{Solution:}
At each step $k$, the number of sets in the partition doubles.
$|P_1| = 2^1 = 2$.
Thus, $|P_6| = 2^6 = 64$.

\hrulefill

\textbf{1.58. Show that the three definitions of a partition are logically equivalent.}
(Proof sketched in conceptual analysis, strictly algebraic verification).
\textbf{Solution:}
(a) \textit{1 implies 2:} Def 1 states subsets are pairwise disjoint and cover $A$. Let $x \in A$. Since the union is $A$, $x$ belongs to at least one subset. Since they are pairwise disjoint, $x$ cannot belong to a second subset. Thus $x$ belongs to exactly one subset (Def 2).
(b) \textit{2 implies 3:} Def 2 states $x$ belongs to exactly one subset. Take two subsets $X, Y \in S$. If they are not disjoint, they share an element $x$. But $x$ belongs to exactly one subset, forcing $X=Y$. Thus, any two subsets are equal or disjoint (Def 3).
(c) \textit{3 implies 1:} Def 3 states any two subsets are equal or disjoint. Since a set $S$ cannot contain identical elements (sets are unique elements), any two distinct subsets in $S$ must be disjoint. Therefore, they are pairwise disjoint (Def 1).

\subsection{Section 1.6: Cartesian Products of Sets}

\textbf{1.59. Let $A = \{x, y, z\}$ and $B = \{x, y\}$. Determine $A \times B$.}

\textbf{Solution:}
\[ A \times B = \{(x,x), (x,y), (y,x), (y,y), (z,x), (z,y)\} \]

\hrulefill

\textbf{1.60. Let $A = \{1, \{1\}, \{\{1\}\}\}$. Determine $A \times A$.}

\textbf{Solution:}
\begin{align*}
A \times A = \{ &(1, 1), (1, \{1\}), (1, \{\{1\}\}), \\
&(\{1\}, 1), (\{1\}, \{1\}), (\{1\}, \{\{1\}\}), \\
&(\{\{1\}\}, 1), (\{\{1\}\}, \{1\}), (\{\{1\}\}, \{\{1\}\}) \}
\end{align*}

\hrulefill

\textbf{1.61. For $A = \{a, b\}$, determine $A \times \mathcal{P}(A)$.}

\textbf{Solution:}
$\mathcal{P}(A) = \{\emptyset, \{a\}, \{b\}, \{a, b\}\}$.
\begin{align*}
A \times \mathcal{P}(A) = \{ &(a, \emptyset), (a, \{a\}), (a, \{b\}), (a, \{a, b\}), \\
&(b, \emptyset), (b, \{a\}), (b, \{b\}), (b, \{a, b\}) \}
\end{align*}

\hrulefill

\textbf{1.62. For $A = \{\emptyset, \{\emptyset\}\}$, determine $A \times \mathcal{P}(A)$.}

\textbf{Solution:}
$\mathcal{P}(A) = \{\emptyset, \{\emptyset\}, \{\{\emptyset\}\}, A\}$.
The product has $2 \times 4 = 8$ pairs, pairing each element of $A$ with each element of $\mathcal{P}(A)$.

\hrulefill

\textbf{1.63. For $A = \{1, 2\}$ and $B = \{\emptyset\}$, determine $A \times B$ and $\mathcal{P}(A) \times \mathcal{P}(B)$.}

\textbf{Solution:}
$A \times B = \{(1, \emptyset), (2, \emptyset)\}$.
$\mathcal{P}(A) = \{\emptyset, \{1\}, \{2\}, \{1, 2\}\}$ and $\mathcal{P}(B) = \{\emptyset, \{\emptyset\}\}$.
Their Cartesian product will contain $4 \times 2 = 8$ ordered pairs.

\hrulefill

\textbf{1.64. Describe the graph of the circle whose equation is $x^2 + y^2 = 4$ as a subset of $\mathbb{R} \times \mathbb{R}$.}

\textbf{Solution:}
\[ S = \{(x, y) \in \mathbb{R} \times \mathbb{R} \mid x^2 + y^2 = 4\} \]

\hrulefill

\textbf{1.65. List the elements of the set $S = \{(x,y) \in \mathbb{Z} \times \mathbb{Z} : |x| + |y| = 3\}$.}

\textbf{Solution:}
Since $x$ and $y$ are integers, their absolute values must sum to 3. The integer pairs $(|x|, |y|)$ can be $(3, 0), (0, 3), (2, 1), (1, 2)$.
Elements: $(3,0), (-3,0), (0,3), (0,-3), (2,1), (-2,1), (2,-1), (-2,-1), (1,2), (-1,2), (1,-2), (-1,-2)$. (Forms a diamond shape).

\hrulefill

\textbf{1.66. For $A = \{1, 2\}$ and $B = \{1\}$, determine $\mathcal{P}(A \times B)$.}

\textbf{Solution:}
$A \times B = \{(1,1), (2,1)\}$.
\[ \mathcal{P}(A \times B) = \{\emptyset, \{(1,1)\}, \{(2,1)\}, \{(1,1), (2,1)\}\} \]

\hrulefill

\textbf{1.67. For $A = \{x \in \mathbb{R} : |x-1| \le 2\}$ and $B = \{y \in \mathbb{R} : |y-4| \le 2\}$, give a geometric description of the points in the $xy$-plane belonging to $A \times B$.}

\textbf{Solution:}
$A = [-1, 3]$ and $B = [2, 6]$.
$A \times B$ forms a solid square in the $xy$-plane with vertices at $(-1, 2), (3, 2), (3, 6),$ and $(-1, 6)$.

\hrulefill

\textbf{1.68. For $A = \{a \in \mathbb{R} : |a| \le 1\}$ and $B = \{b \in \mathbb{R} : |b| = 1\}$, give a geometric description of $(A \times B) \cup (B \times A)$.}

\textbf{Solution:}
$A = [-1, 1]$ and $B = \{-1, 1\}$.
$A \times B$ forms two horizontal line segments at $y=1$ and $y=-1$ from $x=-1$ to $x=1$.
$B \times A$ forms two vertical line segments at $x=1$ and $x=-1$ from $y=-1$ to $y=1$.
The union is the \textbf{boundary of the square} centered at the origin with side length 2.

\hrulefill

\textbf{1.69. Let $A = \{1, 2, 3\}$ and $B = \{2, 3, 4\}$. Determine the following:}
(a) $(A \times B) \cap (B \times A)$ \quad (b) $(A \cap B) \times (B \cap A)$

\textbf{Solution:}
\textbf{Part (a):} For an ordered pair $(x,y)$ to be in the intersection, it must be in $A \times B$ ($x \in A, y \in B$) AND in $B \times A$ ($x \in B, y \in A$). This means $x \in A \cap B$ and $y \in A \cap B$.
$A \cap B = \{2, 3\}$.
Intersection is $\{(2,2), (2,3), (3,2), (3,3)\}$.
\textbf{Part (b):} $(A \cap B) \times (B \cap A) = \{2, 3\} \times \{2, 3\} = \{(2,2), (2,3), (3,2), (3,3)\}$.

\hrulefill

\textbf{1.70. Let $R = \{(1,1), (2,1), (3,2), (3,3), (4,2), (4,3)\}$. Find subsets $A, B, C, D$ of the set $\{1, 2, 3, 4\}$ such that $R = ((A \times B) \cup (C \times D)) - (D \times D)$.}

\textbf{Solution:}
This problem is mathematically impossible as stated. We can provide a rigorous proof by contradiction that no such subsets exist:
\begin{enumerate}
    \item By the definition of set difference, an element in $R$ cannot belong to the subtracted set $D \times D$.
    \item We are given that $(1,1) \in R$ and $(3,3) \in R$.
    \item Therefore, $(1,1) \notin D \times D$ and $(3,3) \notin D \times D$. This implies $1 \notin D$ and $3 \notin D$.
    \item Since $1 \notin D$, no ordered pair in $C \times D$ can have a second coordinate of $1$. Thus, $(1,1) \notin C \times D$. Since $(1,1) \in R$, it must be generated by $A \times B$. Hence, $1 \in A$ and $1 \in B$.
    \item Similarly, since $3 \notin D$, $(3,3) \notin C \times D$. Thus, $(3,3)$ must come from $A \times B$. Hence, $3 \in A$ and $3 \in B$.
    \item Because $1, 3 \in A$ and $1, 3 \in B$, the Cartesian product $A \times B$ \textit{must} contain the cross terms $(1,3)$ and $(3,1)$.
    \item Notice that $(1,3)$ and $(3,1)$ are \textit{not} in $R$. Since they are generated in the union $(A \times B) \cup (C \times D)$, the only way they do not appear in $R$ is if they are subtracted out by $D \times D$.
    \item For $(1,3) \in D \times D$, we must have $1 \in D$ and $3 \in D$. But we already proved in Step 3 that $1 \notin D$ and $3 \notin D$.
\end{enumerate}
This is a fatal contradiction. No such sets exist.

\hrulefill

\textbf{1.71. For $i=1, 2, 3$, let $A_i = \{i, i+1\}$. Determine $\bigcup_{i=1}^3 (A_i \times A_i)$.}

\textbf{Solution:}
$A_1 \times A_1 = \{(1,1), (1,2), (2,1), (2,2)\}$
$A_2 \times A_2 = \{(2,2), (2,3), (3,2), (3,3)\}$
$A_3 \times A_3 = \{(3,3), (3,4), (4,3), (4,4)\}$
The union combines these, merging the overlapping elements $(2,2)$ and $(3,3)$:
\[ \{(1,1), (1,2), (2,1), (2,2), (2,3), (3,2), (3,3), (3,4), (4,3), (4,4)\} \]

\hrulefill

\textbf{1.72. For $A = \{1, 2\}$ and $B = \{3, 4\}$, it follows that $P = \{A, B\}$ is a partition of $S = \{1, 2, 3, 4\}$. Describe a partition of $S \times S$ in terms of Cartesian products of $A$ and $B$.}

\textbf{Solution:}
Since $S = A \cup B$ and $A \cap B = \emptyset$, we can distribute the Cartesian product:
\[ S \times S = (A \cup B) \times (A \cup B) = (A \times A) \cup (A \times B) \cup (B \times A) \cup (B \times B) \]
Because $A$ and $B$ are disjoint, these four Cartesian products are mutually disjoint and non-empty. Thus, the partition is:
\[ \{ A \times A, A \times B, B \times A, B \times B \} \]

\newpage
\section{Part 3: In-Depth Analysis of Problems and Techniques}

\subsection{A) Problem Types and General Approach}
\begin{itemize}
    \item \textbf{Infinite Interval Limits (e.g., 1.45, 1.46):} These problems require you to evaluate the union and intersection of infinite shrinking or expanding intervals. 
    \textit{Strategy:} Evaluate the bounds for $n=1, n=2$ to establish a directional trend. Then evaluate the formal limit as $n \to \infty$. To check if the limit bounds are inclusive or exclusive, test if the limiting value satisfies the condition for *all* $n$ (for intersection) or for *any* $n$ (for union).
    \item \textbf{Partition Verification (e.g., 1.48, 1.49):} 
    \textit{Strategy:} Run the 3-step checklist: 1) Are there any empty sets? 2) Is there any overlap between sets? 3) Does combining all sets recreate the original universal set perfectly? If any answer fails, it is not a partition.
    \item \textbf{Cartesian Geometry (e.g., 1.64, 1.67, 1.68):} Connecting algebraic products to visual planes.
    \textit{Strategy:} Treat the first set as the $x$-axis constraints and the second set as the $y$-axis constraints. 
    \item \textbf{Proof by Exhaustion/Contradiction (e.g., 1.44, 1.70):}
    \textit{Strategy:} For impossible logic puzzles (like 1.70), assume the solution exists, trace the implication of a specific element (like $(1,1)$), and follow the logic until a contradiction occurs.
\end{itemize}

\subsection{B) Key Algebraic and Set-Theoretic Tricks}
\begin{itemize}
    \item \textbf{The Parity Overlap Check (Problem 1.40):} Noticing that $\{i-1, i+1\}$ always produces either two even numbers or two odd numbers. Comparing index $i$ to $i+1$ guarantees disjointness because one will output evens and the other odds.
    \item \textbf{Strictly Nested Sequences (Problem 1.38c):} If $C_1 \supset C_2 \supset C_3 \dots$, then the union of the collection is simply the largest/first set $C_1$. This saves you from having to do infinite algebra.
    \item \textbf{The Cartesian Intersection Identity (Problem 1.69):} Utilizing the rule $(A \times B) \cap (C \times D) = (A \cap C) \times (B \cap D)$.
\end{itemize}

\newpage
\section{Part 4: "Cheatsheet" and Tips for Success}

\subsection{Summary of Crucial Set Rules}
\begin{itemize}
    \item \textbf{Partition Checklist:} 1. No $\emptyset$. \quad 2. Mutually Disjoint ($A_i \cap A_j = \emptyset$). \quad 3. Total Union equals Universe.
    \item \textbf{De Morgan's Laws for Collections:} 
    \[ \overline{\bigcup_{\alpha \in I} A_\alpha} = \bigcap_{\alpha \in I} \overline{A_\alpha} \quad \text{and} \quad \overline{\bigcap_{\alpha \in I} A_\alpha} = \bigcup_{\alpha \in I} \overline{A_\alpha} \]
    \item \textbf{Cartesian Distribution:} $A \times (B \cup C) = (A \times B) \cup (A \times C)$
\end{itemize}

\subsection{Shortcuts and Traps}
\begin{itemize}
    \item \textbf{TRAP: The "Infinity is a Number" Fallacy.} When evaluating intervals like $(-\frac{1}{n}, 2-\frac{1}{n})$ as $n \to \infty$, you must determine if the bound is ever actually reached. The limit bounds an open interval but often creates a half-closed intersection.
    \item \textbf{SHORTCUT: Visualizing Cartesian Products.} If $A$ is finite (size $m$) and $B$ is finite (size $n$), visualize $A \times B$ as an $m \times n$ grid of dots.
    \item \textbf{TRAP: Subsets vs Elements in Partitions.} Remember that $S_4 = A$ is NOT a partition of $A$ (from Problem 1.49). A partition is a set of sets. It must be written as $\{A\}$.
\end{itemize}

\newpage
\section{Part 5: Conceptual Synthesis and The "Big Picture"}

\subsection{A) Thematic Connections}
The core theme of this chapter is \textbf{Organization and Scaling}. We are taking the basic building blocks of set theory (union, intersection, pairs) and scaling them to infinity (indexed collections) and organizing them systematically (partitions). 
This mirrors what you did in Calculus: you learned the area of a rectangle (basic), and scaled it to an infinite sum of infinitesimally thin rectangles (integration).

\subsection{B) Forward and Backward Links}
\textbf{Backward Link:} This chapter relies heavily on solving basic absolute value inequalities (e.g., $|x-1| \le 2$). The algebraic foundation allows you to define the sets before operating on them.
\textbf{Forward Link:} 
\begin{itemize}
    \item \textbf{Topology and Real Analysis:} Understanding open/closed infinite intersections is the direct foundation for Topology. (e.g., an infinite intersection of open sets is not necessarily open, as seen in Problem 1.46a).
    \item \textbf{Probability Theory:} A partition is identical to the concept of "mutually exclusive and exhaustive events" which is required to use the Law of Total Probability or Bayes' Theorem.
\end{itemize}

\newpage
\section{Part 6: Real-World Application and Modeling}

\subsection{A) Concrete Scenarios (Finance \& Stochastic Calculus Focus)}
\begin{enumerate}
    \item \textbf{Derivative Pricing (Information Filtrations):} In stochastic calculus used by quantitative analysts to price derivatives (like European Options), the flow of time and information is modeled using a sequence of partitions called a \textit{filtration}. As time progresses, the "universe" of possible stock paths is partitioned into smaller, more precise subsets reflecting the information gained.
    \item \textbf{Risk Management (State Spaces):} A bank tracking the default risk of two correlated assets sets up a Cartesian product. If Asset 1 has states $A = \{\text{Default, Safe}\}$ and Asset 2 has $B = \{\text{Default, Safe}\}$, the joint probability space is $A \times B$.
\end{enumerate}

\subsection{B) Model Problem Setup: Binomial Tree Filtration}
\textbf{Scenario:} You are pricing an option on a stock over 2 days. The stock can either go Up (U) or Down (D) each day.
\textbf{Variables and Sets:}
Let the universe of paths be $\Omega = \{UU, UD, DU, DD\}$.
Define time $t=0, 1, 2$. At time $t$, we create a partition $\mathcal{F}_t$ of $\Omega$ representing what we know.
\textbf{The Model:}
At $t=0$, we know nothing: $\mathcal{F}_0 = \{\Omega\}$.
At $t=1$, we know the first day's move. The partition splits based on the first letter:
$\mathcal{F}_1 = \{\{UU, UD\}, \{DU, DD\}\}$.
At $t=2$, we know everything. The partition is completely refined:
$\mathcal{F}_2 = \{\{UU\}, \{UD\}, \{DU\}, \{DD\}\}$.
In financial mathematics, the expectation of the payoff is taken iteratively backward through these partitions (Conditional Expectation given a $\sigma$-algebra).

\newpage
\section{Part 7: Common Variations and Untested Concepts}

The homework provided excellent coverage, but omitted a crucial advanced topic frequently tested in university set theory exams: \textbf{Limit Supremum ($\limsup$) and Limit Infimum ($\liminf$) of Sets.}

\textbf{Concept:} Just as sequences of numbers have limits, sequences of sets have limits.
\begin{itemize}
    \item $\liminf_{n \to \infty} A_n = \bigcup_{k=1}^\infty \bigcap_{n=k}^\infty A_n$ (Elements that are in \textit{all but finitely many} sets).
    \item $\limsup_{n \to \infty} A_n = \bigcap_{k=1}^\infty \bigcup_{n=k}^\infty A_n$ (Elements that are in \textit{infinitely many} sets).
\end{itemize}

\textbf{Example Problem:} Let $A_n = [0, 2 - (-1)^n]$. Find $\limsup A_n$ and $\liminf A_n$.
\textbf{Solution:}
If $n$ is even: $A_n = [0, 1]$.
If $n$ is odd: $A_n = [0, 3]$.
The limit infimum requires elements to be in the sets permanently after some point. Since it oscillates, the only permanent elements are those in the smaller interval: $\liminf A_n = [0, 1]$.
The limit supremum requires elements to appear infinitely often. Since the larger interval $[0, 3]$ appears infinitely many times (every odd step), $\limsup A_n = [0, 3]$.

\newpage
\section{Part 8: Advanced Diagnostic Testing: "Find the Flaw"}

Below are five problems with solutions that contain a subtle but critical flaw. Act as the grader: identify the error, explain why it is wrong in one sentence, and provide the correct step.

\textbf{Problem 1:} Let $A_n = (0, \frac{1}{n})$. Find $\bigcap_{n=1}^\infty A_n$.
\textit{Flawed Solution:} As $n \to \infty$, the upper bound $\frac{1}{n}$ approaches $0$. The lower bound is $0$. Therefore, the interval closes in on $0$. The intersection is $\{0\}$.

\textbf{Problem 2:} For sets $A, B, C$, simplify $A \times (B \cap C)$.
\textit{Flawed Solution:} By the distributive property, we can distribute the cross product over the intersection. Therefore, $A \times (B \cap C) = (A \cap B) \times (A \cap C)$.

\textbf{Problem 3:} Determine if $S = \{ \mathbb{R}, \emptyset \}$ is a partition of the real numbers $\mathbb{R}$.
\textit{Flawed Solution:} The union of $\mathbb{R}$ and $\emptyset$ is $\mathbb{R}$. The intersection of $\mathbb{R}$ and $\emptyset$ is $\emptyset$, so they are mutually disjoint. Thus, all conditions are met, and $S$ is a partition.

\textbf{Problem 4:} Let $|A| = 3$ and $|B| = 2$. Find $|\mathcal{P}(A \times B)|$.
\textit{Flawed Solution:} The size of a power set is $2^n$. Because Cartesian product elements are combined, the total power set size is the product of the power sets: $|\mathcal{P}(A \times B)| = 2^3 \times 2^2 = 8 \times 4 = 32$.

\textbf{Problem 5:} Evaluate $\bigcup_{n=1}^\infty [n, n+1]$.
\textit{Flawed Solution:} The sets are $[1, 2], [2, 3], [3, 4]$, etc. They chain together continuously. The starting point is $1$, and as $n \to \infty$, the upper bound goes to $\infty$. Therefore, the union is $[1, \infty]$.

\vspace{1cm}
\hrulefill
\vspace{0.5cm}

\textbf{\Large Grader Answer Key (Your Task):}
\begin{itemize}
    \item \textbf{Problem 1 Flaw:} You cannot include $0$ in the intersection because the interval $(0, \frac{1}{n})$ is open at $0$; since $0 \notin A_1$, it cannot be in the intersection. \textit{Correction:} The intersection is $\emptyset$.
    \item \textbf{Problem 2 Flaw:} The distributive property for Cartesian products distributes over the second sets, keeping the Cartesian operator, not swapping it for an intersection. \textit{Correction:} $A \times (B \cap C) = (A \times B) \cap (A \times C)$.
    \item \textbf{Problem 3 Flaw:} A partition cannot, by strict formal definition, contain the empty set. \textit{Correction:} $S$ is not a partition because it contains $\emptyset$.
    \item \textbf{Problem 4 Flaw:} The cardinality of a Cartesian product is $|A| \times |B| = 6$, so the power set formula is $2^{|A \times B|}$, not the product of the powers. \textit{Correction:} $2^6 = 64$.
    \item \textbf{Problem 5 Flaw:} Infinity is a concept, not a real number, and therefore can never be enclosed by a closed bracket $[ \dots ]$. \textit{Correction:} The union is $[1, \infty)$.
\end{itemize}

\end{document}